\kbvsechead{ಇದು ನಾಮಪದ:}

\indent ಧಾತುಗಳಿಗೆ ಕೃತ್‌ ಪ್ರತ್ಯಯ ಸೇರಿ ಆಗುವ ಶಬ್ದವೇ ಕೃದಂತ, ಇದನ್ನು ಕೃದಂತ ನಾಮಪ್ರಕೃತಿ ಅಥವಾ ಕೃನ್ನಾಮ ಪ್ರಕೃತಿ ಎಂದು ಕರೆಯುವುದು ವಾಡಿಕೆ. ಈ ರೀತಿಯಾಗಿ ನಾಮಪ್ರಕೃತಿಗಳಿಗೆ ನಾಮ ವಿಭಕ್ತಿ ಪ್ರತ್ಯಯಗಳು ಸೇರಿ ನಾಮಪದಗಳಾಗುವುದು; ಇವೇ ಕೃದಂತ ನಾಮಪದಗಳು.\par

\vspace{3pt plus 2pt}\noindent\hspace*{1.2em}\kbvlabel{ಉದಾ}\par\vspace{1pt}

\begin{kbvderiv}
\vspace{4pt plus 3pt minus 1pt}
\begin{center}\begin{adjustbox}{max width=\textwidth}
{\renewcommand{\arraystretch}{1.22}\setlength{\tabcolsep}{4pt}
\begin{tabular}{l c l}
ಓಡುವ + ಅವನು + ಉ & {\color{kbvgrey}$=$} & \kbvresult{ಓಡುವವನು} \\
ಓಡುವ + ಅವನು + ಅನ್ನು & {\color{kbvgrey}$=$} & \kbvresult{ಓಡುವವನನ್ನು} \\
ಓಡುವ + ಅವನು + ಇಂದ & {\color{kbvgrey}$=$} & \kbvresult{ಓಡುವವನಿಂದ} \\
ಓಡುವ + ಅವನು + ಇಗೆ & {\color{kbvgrey}$=$} & \kbvresult{ಓಡುವವನಿಗೆ} \\
ಓಡುವ + ಅವನು + ದೆಸೆಯಿಂದ & {\color{kbvgrey}$=$} & \kbvresult{ಓಡುವವನದೆಸೆಯಿಂದ} \\
ಓಡುವ + ಅವನು + ಅ & {\color{kbvgrey}$=$} & \kbvresult{ಓಡುವವನ} \\
ಓಡುವ + ಅವನು + ಅಲ್ಲಿ & {\color{kbvgrey}$=$} & \kbvresult{ಓಡುವವನಲ್ಲಿ} \\
ಓಡುವ + ಅವನು + ಏ & {\color{kbvgrey}$=$} & \kbvresult{ಓಡುವವನೇ} \\
\end{tabular}}
\end{adjustbox}\end{center}
\vspace{4pt plus 3pt minus 1pt}
\end{kbvderiv}

\indent ಮೇಲಿನ ಉದಾಹರಣೆಯಲ್ಲಿ ಕೃದಂತ ನಾಮಪ್ರಕೃತಿಗೂ ನಾಮ ವಿಭಕ್ತಿ ಪ್ರತ್ಯಯಕ್ಕೂ ನಡುವೆ “ಅವನು” ಎಂಬ ಸರ್ವನಾಮ ಆಗಮವಾದುದು. ಇದೇ ರೀತಿಯಾಗಿ, ಸ್ತ್ರೀಲಿಂಗ ಏಕವಚನದಲ್ಲಿ “ಅವಳು” ಮತ್ತು ನಪುಂಸಕಲಿಂಗ ಏಕವಚನದಲ್ಲಿ “ಅದು” ಎಂಬ ಸರ್ವನಾಮಗಳು ಬರುವುದು.\par

\begin{kbvrule}
\noindent\kbvbullet\hspace{0.5em}ಇದರ ಸಹಾಯದಿಂದ ನಾಮಪದಗಳನ್ನು ವಿಶೇಷಿಸುತ್ತಾರೆ.
\end{kbvrule}

\vspace{3pt plus 2pt}\noindent\hspace*{1.2em}\kbvlabel{ಉದಾ}\par\vspace{1pt}

\noindent\hspace*{1.8em}ಮಾ ಅಸ... (ಭೂತ ಕೃದಂತ)\hspace{0.9em}ಹಾ ಬಾರಿ... (ಭವಿಷ್ಯತ್‌ ಕೃದಂತ)\hspace{0.9em}ಸೂಚಕೆ; ಓಡ (ನಿಷೇಧ ಕೃದಂತ)\par\vspace{3pt plus 2pt}

\noindent\hspace*{2.2em}ನಿಷೇಧಾರ್ಥ ಸೂಚನೆ: “ಓಡದ” ಇಲ್ಲಿ “ಅದ” ಎಂಬುದು.\par
